\chapter{Ectopic-Pregnancy Surgery}

\section*{Introduction \& Clinical Indications}
An ectopic pregnancy implants outside the uterine cavity, most commonly in a fallopian tube. Surgery is required when rupture or hemodynamic instability is suspected, when medical treatment is unsuitable or unsuccessful, or when disease characteristics and follow-up circumstances make an operation the safest option. A stable patient may instead be eligible for methotrexate or expectant management after specialist assessment.

Ruptured ectopic pregnancy is a time-critical emergency. Resuscitation, blood preparation, gynecologic consultation, and operative control of bleeding proceed in parallel. The operation does not treat a pregnancy of unknown location until the diagnosis and clinical situation have been evaluated appropriately.

\section*{Relevant Surgical Anatomy}
The fallopian tube lies between the ovary and uterine horn and is closely related to the uterine and ovarian vessels, broad-ligament structures, bowel, bladder, and ureter. The ampulla is the commonest tubal implantation site. Contralateral tubal condition, ovarian reserve, prior surgery, and the patient’s fertility goals influence whether tube-preserving treatment is appropriate.

\section*{Preoperative Workup \& Preparation}
Assess symptoms, vital signs, serial quantitative hCG, transvaginal ultrasound, hemoglobin, blood group and antibody status, renal and liver function, and suitability for anesthesia or methotrexate. Determine whether rupture is likely and review the opposite tube, fertility plans, prior ectopic pregnancy, and ability to return for follow-up. Explain laparoscopy, possible laparotomy, salpingectomy, salpingostomy, conversion, transfusion, persistent trophoblast, recurrence, and future fertility implications.

\section*{Surgical Technique \& Procedure}
Stable patients commonly undergo laparoscopy. The pelvis is inspected, bleeding controlled, and the affected tube removed or opened to remove the pregnancy according to tubal damage, contralateral anatomy, fertility goals, and disease factors. A salpingostomy requires postoperative hCG surveillance for persistent trophoblastic tissue. Unstable patients or those with major hemoperitoneum may require immediate open surgery and blood-product support.

\section*{Complications \& Risks}
Risks include hemorrhage, transfusion, infection, injury to bowel, bladder, ureter, or vessels, anesthesia complications, adhesions, persistent trophoblast, recurrent ectopic pregnancy, and reduced fertility. Shoulder-tip pain and emotional distress may occur during recovery. New severe pain, dizziness, fainting, fever, heavy bleeding, or shoulder pain with weakness requires urgent evaluation.

\section*{Postoperative Care \& Recovery}
Monitor vital signs, hemoglobin, pain, wound sites, and hCG when salpingostomy or incomplete resolution is possible. Provide anti-D prophylaxis when indicated by local guidance and blood group, contraception or fertility counseling, and emotional support. Future pregnancies need early ultrasound to confirm location because recurrence risk is increased.

\section*{Outcomes \& Prognosis}
Most early ectopic pregnancies can be treated successfully with surgery or medical therapy. Fertility after surgery depends on the condition of the remaining tube, pelvic disease, age, and prior fertility. Early diagnosis and rapid treatment of rupture are critical to preventing avoidable morbidity and death.

\section*{References}
\begin{enumerate}
\item American College of Obstetricians and Gynecologists. Practice Bulletin: Tubal Ectopic Pregnancy.
\item Royal College of Obstetricians and Gynaecologists. Diagnosis and Management of Ectopic Pregnancy.
\item National Institute for Health and Care Excellence. Ectopic pregnancy and miscarriage: diagnosis and initial management.
\end{enumerate}
